\subsection{The normalized pairing of two digit symbols}
\label{sec:block-pairing}
\hypertarget{source-block-pairing}{}

Ranks count the retained symbols, but combining the segment tests at an edge also requires their pairings. The two models use factorials at different degrees. Here we
compare their pairings before lifting off the symbol. All
valuations in this calculation have \(v(p)=1\) and \(w^2=p\).

\paragraph{Factorial units.}

To compare the coarse and true factorial pairings, separate the unit part of the factorial of an exponent with prescribed low digit.

For \(P=p^d\), \(d\ge1\), and \(0\le h<P\), write
\[
 (h+Pn)!=h!\,n!\,(P!)^n C_P(n)D_{P,h}(n),\qquad
 D_{P,h}(n)=\prod_{1\le j\le h}(1+Pn/j).
\]
Here \(C_P(0)=1\) and
\(C_P(n+1)/C_P(n)=\prod_{1\le j<P}(1+Pn/j)\),
as is seen by cancelling the last multiple of \(P\) in the
factorial quotient. Every \(P/j\) in these products belongs
to \(p\mathbb Z_p\). Both factors are units.

\begin{lemma}[Approximating inverse factorial units]
\label{lem:factorial-unit-polynomial}
For any finite product of the factors \(C_P(\mu_i)^{\pm1}\) and
\(D_{P,h}(\mu_i)^{\pm1}\), with the indicated restrictions on \(P,h\),
and any positive integer \(c\), write \(A(\mu)\) for that product. One polynomial \(P_c\) works for every nonnegative integral vector \(\mu\),
\[
 \deg P_c\le4c,\qquad P_c(\mu)\equiv A(\mu)\pmod{p^c}.
\]
 Its coefficients in the product
binomial basis belong to \(\mathbb Q\cap\mathbb Z_p\).
\end{lemma}

\begin{proof}
\proofstep{Polynomial interpolation of the factorial units.}
We need a uniform polynomial for the inverse units modulo
\(p^c\), with degree controlled independently of the radices,
digits and number of variables. We include a finite difference
proof. In each product \(B(x)=\prod_{j\le h}(1+Px/j)^{\pm1}\),
the coefficient of \(x^e\) has valuation at least \(e\). This
also holds for negative exponents by the convergent geometric
series, uniformly on \(\mathbb Z_p\). For the forward difference
\(\nabla f(x)=f(x+1)-f(x)\),
\[
 v(\nabla^l(B-1)(x))\ge\max(1,l)
       \quad (l\ge0,\ x\in\mathbb Z_{\ge0}).
\]
Indeed the constant term of \(B-1\) is zero and the \(l\)-th
difference kills powers of degree less than \(l\).

For \(F=C_P^{\pm1}\) we have \(\nabla F=(B-1)F\) with \(h=P-1\).
The shifted product rule gives
\[
 \nabla^j F(x)
   =\sum_{l=0}^{j-1}\binom{j-1}{l}
       \nabla^l(B-1)(x)\,\nabla^{j-1-l}F(x+l).
\]
Using the bound on \(B-1\) and induction on \(j\), each term
has valuation at least \(\max(1,l)+(j-1-l)/2\ge j/2\). The same bound holds for
\(D_{P,h}^{\pm1}\) by the coefficient estimate. Repeated use of
the product rule preserves this bound for products in one
variable.

For a product of such functions in several variables
the mixed difference of order \(\nu\) at zero has valuation
at least \(|\nu|/2\). Newton's finite difference formula at a
nonnegative integer vector \(\mu\) is exact with indices
\(\nu\le\mu\), by binomial inversion in each coordinate.
Truncating to \(|\nu|\le4c\) thus gives a polynomial
\[
 H(\mu)=\sum_{|\nu|\le4c} q_\nu\prod_v\binom{\mu_v}{\nu_v},
 \qquad q_\nu\in\mathbb Q\cap\mathbb Z_p,
\]
congruent to the desired product modulo \(p^c\) for every
\(\mu\). A variable on which the unit is constant occurs with
degree zero. The decomposition applies to any finite product of digit factors.

\end{proof}

\begin{established}[Injectivity of the block monomial embedding]\label{hyp:digit-block-embedding}
Embed the block variables in the \(k\) coordinates of the
two letters at each interior top and at the next even top.
The \(r\) ordinary coordinates occupy columns 1 through \(r\).
At a tie the extra coordinates \(Z,S\) occupy column zero in
the vector and covector, respectively, including at the final
pair. A block monomial \(X^\mu\) maps to \(X^{I(\mu)}\) by
multiplying its exponents at interior top \(i\) by \(P_i\),
then adding the digit vectors. At an untied interior top also
add \(P_i(A_i-w_i),P_i h_i\) in the two columns zero.
There are no free variables in those two columns.
Writing \(\mathcal M_B\) for the monomials in block \(B\), the map satisfies
\[
 I_B(\mu)=I_B(\nu)\Longrightarrow\mu=\nu,\qquad
 I_B(\mathcal M_B)\cap I_{B'}(\mathcal M_{B'})=\varnothing\quad(B\ne B').
\]
\end{established}
\begin{proof}
At the
end the radix is one and the digits are zero.  On each ordinary coordinate the recovery is Euclidean division. If \(m\) is its true exponent at top \(i\), then
\[
 m=P_i q+\delta,\qquad q=\lfloor m/P_i\rfloor,\quad \delta=m\bmod P_i.
\]
At an untied top the
quotients in column zero recover \(w_i,h_i\). All remaining
coarse exponents are then recovered by division by the radices.
At any tie the difference of the exponents of \(Z\) and \(S\)
is its decoration \(z\), because each power of
\(E=ZS+x\cdot y\) has weight zero for that difference.
This recovers \(o_i=A_i-z_i\) or \(o_+=a_+-z_+\).
Thus different blocks have disjoint monomial supports and
\(I\) is injective within a block.
\end{proof}

\begin{lemma}[The scale of the digit block]\label{lem:digit-block-scale}
For the embedding of Lemma~\ref{hyp:digit-block-embedding}, the factor
\(\Phi\) in \(I(\mu)!=\Phi\,\mu!\,U(\mu)\) is independent
of \(\mu\). For the scale \(\gamma_I\) of \eqref{eq:digit-block-scale} one has
\[
 0\le v(\gamma_I)\le(k+1)L_I.
\]
 These bounds use the
same valuation \(v(p)=1\) as the digit pairing.
\end{lemma}

\begin{proof}
\proofstep{The degrees and the common scalar factor.}
Write \(g_x\) for the sum of the exponents in the free
variables belonging to a particular true letter \(x\), before
inserting digits. It is constant on all the levels of the
block. At an untied pair these sums are \(w_i,C_i-h_i\), at a
tie they are \(A_i,C_i\), and at the end they are \(a_+,b_+\).
Indeed each term of \(E\) has bidegree \((1,1)\), and each
\(Z^z\) restores its prescribed decoration. Write \(D_x\) for
the true degree after adding digits and the two fixed columns.
Put
\[
 b_*=\prod_{i<\sigma}(A_i-w_i)!\,h_i!,\qquad
 T=b_*\prod_x g_x!,\qquad V=v_p(T),\qquad
 L_I=\sum_x\lfloor\log_p D_x\rfloor .
\]
Use a zero log for a zero degree. We have \(b_*\mu!\mid T\)
on this block, by multiplying the letter multinomials.

Let \(U\) be the product of the \(C_PD_{P,h}\) over the free
coordinates at the odd top indices.
The factorial-unit identity gives, for a fixed integer
\(\Phi>0\) depending on this block and all its ambient degrees,
\[
 I(\mu)!=\Phi\,\mu!\,U(\mu).
\]
The factors independent of \(\mu\)
in \(\Phi\) are the factorials of the two fixed column
exponents at each untied top and the digit factorials in the
free coordinates. The remaining factor is
\(\prod_x(P_x!)^{g_x}\), independent of \(\mu\) by the preceding
homogeneity.

Write \(\partial^{[\nu]}\) for the divided derivative.
Apply the polynomial just proved to \(U^{-1}\),
with variables restricted to the odd tops. It gives the
diagonal operator
\[
 \mathcal H=\sum_\nu q_\nu X^\nu\partial^{[\nu]},
 \qquad \mathcal H X^\mu=H(\mu)X^\mu .
\]
Define
\begin{equation}\label{eq:digit-block-scale}
 Z_I=\frac{p^{kL_I}}{\prod_x D_x!},\qquad
 \gamma_I=Z_I p^V\Phi/b_* .
\end{equation}
Then
\[
 0\le v(\gamma_I)\le(k+1)L_I.
\]
To see the
cancellation in this assertion, form an exponent \(\mu^0\)
by putting each undecorated letter degree in its first ordinary
coordinate, all decorations in \(Z\) and zero in \(S\).
This uses level zero and has the right letter degrees,
irrespective of whether it obeys the prefix cuts.

In each
letter \(\mu^0\) has at most two nonzero entries. The
integer \(\prod_xg_x!/\mu^0!\) is a product of binomials,
of valuation between zero and \(L_I\) by the carry bound.
The true multinomials in \(\prod_xD_x!/I(\mu^0)!\) have
valuations summing to at most \((k-1)L_I\).
Legendre's formula proves both bounds, since at each positive
power of \(p\) the difference of the floor sums is at most one
for two entries and \(k-1\) for \(k\) entries. As
\[
 v(I(\mu^0)!)=v(\Phi)+v(\mu^0!),
\]
subtraction gives
\[
 v(\gamma_I)=kL_I+
 v\!\left(\frac{\prod g_x!}{\mu^0!}\right)
 -v\!\left(\frac{\prod D_x!}{I(\mu^0)!}\right).
\]
Both factorial quotients in this formula have nonnegative
valuation. Their upper bounds from the selected monomial
are \(L_I\) and \(kL_I\), respectively. Substituting them
gives the two endpoints for \(v(\gamma_I)\); in particular
\(\gamma_I\) is nonzero and integral.

\end{proof}

\paragraph{The two symbols.}
For two inputs \(u,u'\) in the factorial lattice of the block
take \(s(u)=w^V\mathcal H u\), \(s'(u')=w^V u'\).
As \(b_*\mu!\mid T\), the defining monomial inequality of
that lattice makes \(w^V u\) coefficient-integral. The same
holds after applying \(\mathcal H\).

We verify membership in the symbol lattice at the required
margin. First the integral span of the generators at a fixed
margin is saturated inside its rational span and the
coefficient lattice. Write a rational combination in levels
as a sum of \(\prod E_i^{j_i}Z_i^{z_i}\) times polynomials in
the reduced ordinary letters and the residual \(\xi\)'s.
Set every \(Z_i=1\) and then substitute
\[
 S_i\mapsto S_i-x_i\cdot y_i,
\]
including the last pair. These integral substitutions
take \(E_i\) to \(S_i\). Extracting the coefficients of
\(\prod S_i^{j_i}\xi^\alpha\) recovers separately the
constrained ordinary-letter polynomials. If the original
combination was integral these polynomials belong to the
coefficient lattice of their constrained rational spaces,
and reassembling them gives integral generators.

\begin{proposition}[The two lifts reproduce the block pairing]
\label{prop:digit-symbol-pairing}
Write \(B(u,u')\) for the normalized diagonal Fock pairing of the two block symbols. Then
\[
 B(u,u')-\gamma_I b_*\langle u,u'\rangle\in\gamma_I w^c O.
\]
 This holds for inputs in the two coefficient
lattices, including the allowed denominators. If the generators satisfy
the raw-tail assertion \eqref{eq:raw-tail-required}, the simultaneous
lifts preserve this congruence and give the two integral tests in
\eqref{eq:two-letter-tensor-rank}.
\end{proposition}

\begin{proof}
\proofstep{Compatibility with a variation of one coordinate.}
Next, replacing one coordinate \(v\) by \((1+t)v\) and taking
\([t^b]\) realizes \(v^b\partial_v^{[b]}\).
On the constrained letter part it increases each allowed first
row by at most \(b\). In fact in a prefix frame change a
rescaling of one coordinate becomes a change by \(t\) times a
rank-one endomorphism on that letter. Its degree-\(a\) term
is a sum of polynomials with \(a\) variables replaced and
\(a\) derivatives.

In any conjugate frame, replacing one
variable of a given vector (or covector) letter by another
changes the first torus-coordinate weight by at most one.
This proves the prefix-weight bound for the whole orbit, hence
the highest-weight cutoff. The integral divided-derivative
formula preserves integrality. Rescaling an odd coordinate also
changes at most one \(E_i\) and its adjacent \(Z_i^{z_i}\).
If a term of \(E_i\) is hit, write
\[
 E_i\mapsto E_i+tM.
\]
For an ordinary variable \(M=x_h y_h\).
For \(Z_i\) or \(S_i\),
\[
 M=Z_iS_i=E_i-x_i\cdot y_i.
\]

After expanding, removing \(l\) powers of \(E_i\) inserts
\(l\) ordinary bilinears. Each has vector and covector degrees
one and first-coordinate weight at most one in every prefix
frame (also at the prefix between its letters). Moving from
level \(j_i\) to \(j_i-l\) supplies precisely those degrees.
The power of \(Z_i\) contributes just an integer binomial
coefficient. Thus \([t^b]\) relaxes every cutoff by at most
\(b\) and preserves the span of the integral generators at the
relaxed cutoffs. Iteration costs at most \(|\nu|\le4c\) in
every cutoff. Since
\[
 m_0(C_m)-m_0(13)=(C_m-13)kN_1Qc\ge4c
\]
for \(C_m\ge17\),
both symbols belong to the required integral span at
\(m_0(13)\).

\proofstep{The diagonal Fock pairing.}
The Fock form on monomials has diagonal \(I(\mu)!\).
For all inputs in one block, including those with denominators,
we now have
\[
 Z_I\langle I(s(u)), I(s'(u'))\rangle
   \equiv \gamma_I b_*\langle u,u'\rangle
                  \pmod{\gamma_I w^c O}.
\]
Indeed \(H(\mu)U(\mu)\equiv1\pmod{p^c}\), the two factors of
\(w^V\) give \(p^V\), and each
\(b_*\mu!\,u_\mu u'_\mu\) is integral by the two defining
coefficient inequalities. Distinct blocks pair to zero by
disjoint support. Products of the symbol maps over the segments
multiply the pairings since their true letters are disjoint,
and \(Z_{\mathrm{ten}}=\prod Z_I\). The diagonal multiplier
has depth at most
\[
 (k+1)\sum_I L_I\le4\kappa_0,
\]
since the \(2N\) degrees
are at most \(2a\), \(\log_p(2a)\le2\log_p a\), and
\(\kappa_0=2Nk\log_p a\). This gives the segment-to-symbol
identity, including its uniform relative congruence.

To use the identity on two lifted polynomials it suffices that
their normalized edge substitutions be coefficient-integral
and congruent to the respective product symbols modulo \(w^c\).
On the true letter degrees the monomial diagonal of
\(Z_{\mathrm{ten}}\langle\ ,\ \rangle\) is integral, by the
same multinomial carry bound. Bilinearity then preserves the
congruence modulo \(w^c\). \paragraph{Orders, corrections and duality.}
Sections~\ref{sec:column-lifts} and~\ref{sec:symbol-correction}
construct the lifts used here. The order of a
block at a segment and the order of a tuple are
\[
 o_I=o_++\sum_i\left(P_i u_i+\sum_{h=1}^r\rho_{ih}+\tau_{i0}\right),
 \qquad o=\sum_I o_I .
\]
Let \(S_I=\sum_i\min(b_i,a_{\rm start}/p^{d_i})\). Write
\(R=a/p^2+b'/p+\mathcal L a/(pk)+\mathcal L b'/k+kQ\).
Then \(o_I\le S_I+30(s+1)R\). Indeed the digits in each summand
cost at most \(kQ\), and
\(P_i K_i\le a/p^{d_i}+2b'/p\) since \(P_i B\le a/p^{d_i}\)
and \(P_i\sum_{l<i}A_l\le (4/p^2)P_i A_i\le8b'/p^2\).

In the initial list the minimum in \(S_I\) is \(b_i\) by
the excluded band. The retained upper bound for \(u_i\) gives
\[
 P_iu_i\le b_i(1+4/p^2)+(\mathcal L/k)P_i K_i.
\]
In the tied list the minimum is at least \(a/p^{d_i}\), and
\(P_i u_i\le P_i K_i+2\mathcal L b'/k\).
For the end, put
\[
 C_{\rm end}=a/p^2+3b'/p+
       (\mathcal L/k+4/p^2)4a/p+2\mathcal L b'/k .
\]
The increasing distances are odd, \(P_i\ge p\), and \(G\ge2\).
Hence \(B\le a/p^2\), \(\sum A_i\le3b'/p\) and
\[
 \sum 2a/P_i\le4a/p
\]
by summing the geometric series. The
last retained cutoff implies \(o_+\le C_{\rm end}\). These bounds
give the claimed constant for \(p\ge5\). Summing them and using
\(\sum(s+1)\le N\) proves \(o\le\lfloor t_0\rfloor\)
for \(C_t\ge30\), as the additional \(Qc\) term is nonnegative.

\proofstep{The two lift maps in a basis adapted to the edge.}
Fix an edge line and use its adapted frame as the identity.
Denote by \(S_g\) the normalized substitution on polynomials
in the frame \(g\). The residue lines parametrize neighbors of
the source. Write \(A_1=13kN_1Q\) and \(M_0=12c\).
At precision \(0\le n\le c\) take the integral span of
product symbols with order at most \(o\), prefix margin
\(M_0+nA_1\), and excess
\[
 e=(c-n)A_1
\]
in their block bounds.
The relaxed bounds keep digit compatibility, the unused labels
equal to zero and the usual integer degree ranges. They require
\[
 h_i\le2\mathcal L B/k+e\ (i<\sigma),\qquad
 o_i\le (1-\Delta/4)A_i+e\ (i\ge\sigma),\qquad
 o_+\le\min(a_+,C_{\rm end}+e).
\]
Every retained block satisfies these bounds at \(n=c\).

For
the middle inequality use \(K_i\le B+\delta A_i\), the second
part of \eqref{eq:coarse-split}, and
\(\delta+\mathcal L/k\le3\Delta/4\).
Call this span \(\mathcal S(n,o)\). A polynomial is clean at
order \(o\) if it has the true multidegrees, its product by
\(w^o\) has integral coefficients in the source frame, and
its substitutions at every other residue line belong to
\(w^c O[X]\).

\proofstep{The required raw polynomial tail.}
The raw polynomial assertion is the following. For each
product generator \(b\) of \(\mathcal S(n,o)\), \(1\le n\le c\),
of actual order \(o_b\le o\), there is a clean \(F_0\)
at that order with
\begin{equation}\label{eq:raw-tail-required}
 S_1F_0-b\equiv\sum_{\epsilon=1}^{n-1}w^\epsilon b_\epsilon
 \pmod {w^nO[X]},\qquad
 b_\epsilon\in\mathcal S(n-\epsilon,o_b+\epsilon).
\end{equation}
Given this, induction lifts the entire span exactly to the
required precision. At zero use the zero polynomial. At \(n\)
lift \(b_\epsilon\) by induction and subtract \(w^\epsilon\)
times its lift. The multiplier recovers the lost source order
\(o_b\) and preserves cleanliness at every other line. It
raises the precision of the correction to \(n\). Integral
linear combinations finish the induction. At \(n=c\) the
margin is \(m_0(13)\), precisely the one obtained above, and
we can do this independently on an edge and on its dual.

\proofstep{The covector in the integral dual.}
Here is the identification of the second lift with a covector.
In polynomial coordinates let
\[
 \Xi(y)(x)=Z_{\rm ten}\langle x,y\rangle.
\]
Multihomogeneous polynomial coordinates realize
the tensor of one-row vector and covector representations,
with the monomial coefficient lattices in the corresponding
frames. The Fock form transports between a frame and its
inverse transpose. Indeed it is obtained by contracting
all indices of two symmetric tensors, summed over the
permutations in each letter, or equivalently by evaluating
the constant term of \(P(\partial)Q\). This works for both
letter orientations. On integral coefficients its normalized
form is integral in any such paired frames.

At a vertex
with frame \(b\), the dual frame is \(b^{-t}\). Along an edge
write \(D(u)=\operatorname{diag}(u,1,\ldots,1)\).
The intermediate frame is \(bD(w^{-1})\).
On the dual edge starting at the dual of the original target
the intermediate frame agrees with the inverse transpose, since
\(w^2=p\) and \(D\) is diagonal:
\[
 b^{-t}D(p)D(w^{-1})=b^{-t}D(w)=(bD(w^{-1}))^{-t}.
\] This remains true on all
Frobenius-twisted coefficient frames, since the twist on
\(b\) fixes \(p\). Thus \(\Xi\) takes the vertex and intermediate
integral lattices on the dual edge into the duals of those
on the original edge. In particular for a lattice \(L\) and an integer \(a\), linearity gives
\[
 \Xi(w^a L)=w^a\Xi(L).
\]

In particular the two normalized intermediate lattices both
carry \(w^{-A_H}\), where \(A_H\) is the sum of vector
degrees. The order multiplier \(w^{-t}\) at the vertex and the
conductor multiplier \(w^c\) at every other intermediate
lattice transport with the same exponents. Duality bijects
outgoing neighbors of the dual source with incoming neighbors
of the original target, preserving the distinguished pair.
This proves the target membership required of the second lift.
In the distinguished frames the normalized substitution of
each polynomial carries \(w^{A_H}\) times its expression in
its intermediate frame. Hence pairing its two substitutions
with \(Z_{\rm ten}\) agrees with evaluation by \(\Xi\) after applying the edge normalization
\[
 w^{A_H}\cdot w^{A_H}=w^{2A_H}=p^{A_H}.
\]
This gives the congruence used in
\eqref{eq:two-letter-tensor-rank} using
\eqref{eq:raw-tail-required}, proved in Section~\ref{sec:symbol-correction}.
\end{proof}
